\documentclass[a4paper,11pt]{article}
\usepackage[utf8]{inputenc}
\usepackage[spanish]{babel}
\usepackage{geometry}
\usepackage{enumitem}
\usepackage{sectsty}
\usepackage{titlesec}
\usepackage{tikz}
\usepackage{array}

\geometry{
    a4paper,
    left=20mm,
    top=18mm,
    right=20mm,
    bottom=20mm
}

% Sin números de página
\pagestyle{empty}

% Estilo de secciones con números en círculos
\newcommand{\circled}[1]{%
    \tikz[baseline=(char.base)]{%
        \node[shape=circle,draw,inner sep=0.5pt,fill=black,text=white,font=\bfseries\footnotesize] (char) {#1};%
    }%
}

% Formato de secciones
\titleformat{\section}
{\normalfont\Large\bfseries}
{\circled{\thesection}}{0.5em}{}

\titlespacing*{\section}{0pt}{12pt}{8pt}

\titleformat{\subsection}
{\normalfont\normalsize\bfseries}
{}{0em}{}

\titlespacing*{\subsection}{8pt}{4pt}{4pt}

\begin{document}

% Encabezado del curso
\begin{center}
\small{Universidad de Granada}

\vspace{0.3cm}

\textbf{\large{Producción e Interacción Oral - Nivel 7}}

\vspace{0.4cm}

\small
\begin{tabular}{@{}ll@{}}
\textbf{Grupos:} & A (Aula 24) y B (Aula 08) \\
\textbf{Centro:} & Centro de Lenguas Modernas / Edificio A \\
\textbf{Profesor:} & Javier Benítez Lainez \\
\textbf{Tutorías:} & Martes 9:00-10:00 y 14:30-15:30 \\
& Sala de profesores / Email: benitezl@go.ugr.es
\end{tabular}
\end{center}

\vspace{0.5cm}
\section*{Repaso Post-Vacaciones}

\vspace{0.4cm}

\section{Reactivación de Conocimiento y Habilidades C1}

\vspace{0.2cm}

---

\vspace{0.2cm}

\section{INTRODUCCIÓN}

\vspace{0.2cm}

Este documento ayuda a reactivar conocimientos y habilidades después de un período vacacional, asegurando que mantengas tu nivel C1 de español oral.

\vspace{0.2cm}

---

\vspace{0.2cm}

\section{DIAGNÓSTICO POST-VACACIONES}

\vspace{0.2cm}

\subsection{Autoevaluación rápida}

\vspace{0.2cm}

Responde estas preguntas honestamente:

\vspace{0.2cm}

\textbf{Gramática:}

1. ¿Recuerdas cuándo usar ser vs. estar? ( ) Sí ( ) Más o menos ( ) No

2. ¿Recuerdas cuándo usar subjuntivo? ( ) Sí ( ) Más o menos ( ) No

3. ¿Recuerdas la diferencia entre pretérito indefinido e imperfecto? ( ) Sí ( ) Más o menos ( ) No

4. ¿Recuerdas los tipos de condicionales? ( ) Sí ( ) Más o menos ( ) No

\vspace{0.2cm}

\textbf{Vocabulario:}

5. ¿Puedes recordar 5 conectores diferentes? ( ) Sí ( ) No

6. ¿Recuerdas vocabulario específico de temas trabajados? ( ) Sí ( ) Más o menos ( ) No

\vspace{0.2cm}

\textbf{Fluidez:}

7. ¿Cómo se siente tu fluidez? ( ) Igual ( ) Un poco oxidada/o ( ) Muy oxidada/o

\vspace{0.2cm}

---

\vspace{0.2cm}

\section{REPASO INTENSIVO (3 DÍAS)}

\vspace{0.2cm}

\subsection{DÍA 1: GRAMÁTICA CLAVE}

\vspace{0.2cm}

\section{Sesión 1 (30 min): Ser vs. Estar}

\vspace{0.2cm}

\textbf{Regla rápida:}

\begin{itemize}

  \item \textbf{SER} = identidad, profesión, características permanentes, origen

  \item \textbf{ESTAR} = estados, emociones, ubicaciones, condiciones temporales

\end{itemize}

\vspace{0.2cm}

\textbf{Práctica:}

Completa con ser o estar:

1. "Yo ____ de México." → Soy

2. "____ muy cansado hoy." → Estoy

3. "____ médico." → Soy

4. "La casa ____ muy grande." → Es

5. "____ en el parque." → Estoy

\vspace{0.2cm}

---

\vspace{0.2cm}

\section{Sesión 2 (30 min): Subjuntivo}

\vspace{0.2cm}

\textbf{Regla rápida:}

Usa subjuntivo después de:

\begin{itemize}

  \item Dudas: "No creo que..."

  \item Deseos: "Quiero que..."

  \item Frases impersonales: "Es importante que..."

\end{itemize}

\vspace{0.2cm}

\textbf{Práctica:}

Transforma a subjuntivo:

1. "Creo que \textbf{viene}." → "No creo que \textbf{venga}."

2. "Es seguro que \textbf{es} verdad." → "Es posible que \textbf{sea} verdad."

3. "Quiero que \textbf{tienes} razón." → "Quiero que \textbf{tengas} razón."

\vspace{0.2cm}

---

\vspace{0.2cm}

\section{Sesión 3 (30 min): Pretéritos}

\vspace{0.2cm}

\textbf{Regla rápida:}

\begin{itemize}

  \item \textbf{Indefinido}: evento específico completado

  \item \textbf{Imperfecto}: hábito, descripción, estado en pasado

  \item \textbf{Pluscuamperfecto}: anterior a otro evento pasado

\end{itemize}

\vspace{0.2cm}

\textbf{Práctica:}

Elige el pretérito correcto:

1. "Ayer ____ (fui/iba) al cine." → Fui

2. "____ (Era/Fue) una noche oscura y ____ (llovía/llovió)." → Era, llovía

3. "Ya ____ (comí/había comido) cuando llegaste." → había comido

\vspace{0.2cm}

---

\vspace{0.2cm}

\subsection{DÍA 2: VOCABULARIO Y CONECTORES}

\vspace{0.2cm}

\section{Sesión 1 (30 min): Conectores}

\vspace{0.2cm}

\textbf{Repaso rápido de conectores:}

\vspace{0.2cm}

\textbf{Organización:}

\begin{itemize}

  \item Primero / En primer lugar / Para empezar

  \item Luego / Después / A continuación

  \item Por último / Finalmente

\end{itemize}

\vspace{0.2cm}

\textbf{Contraste:}

\begin{itemize}

  \item Sin embargo / No obstante / Por el contrario

\end{itemize}

\vspace{0.2cm}

\textbf{Adición:}

\begin{itemize}

  \item Además / Asimismo / Por otro lado

\end{itemize}

\vspace{0.2cm}

\textbf{Causa-Consecuencia:}

\begin{itemize}

  \item Por lo tanto / Por consiguiente / En consecuencia

\end{itemize}

\vspace{0.2cm}

\textbf{Práctica:}

Usa 3 conectores diferentes en una oración: "____ (Primero), hice ejercicio. ____ (Después), fui a casa. ____ (Finalmente), cené."

\vspace{0.2cm}

---

\vspace{0.2cm}

\section{Sesión 2 (30 min): Vocabulario Temático}

\vspace{0.2cm}

\textbf{Repaso de temas clave:}

\vspace{0.2cm}

\textbf{Opinión y Argumentación:}

\begin{itemize}

  \item En mi opinión / Considero que / Me parece que

\end{itemize}

\vspace{0.2cm}

\textbf{Acuerdo/Desacuerdo:}

\begin{itemize}

  \item Estoy de acuerdo / No puedo estar más de acuerdo

  \item Disiento / No estoy convencido

\end{itemize}

\vspace{0.2cm}

\textbf{Tiempo:}

\begin{itemize}

  \item Anteriormente / Posteriormente / Simultáneamente

  \item A lo largo de / Durante / Mientras

\end{itemize}

\vspace{0.2cm}

\textbf{Práctica:}

Escribe 3 oraciones sobre tu última experiencia de vacaciones usando vocabulario de opinión y tiempo.

\vspace{0.2cm}

---

\vspace{0.2cm}

\subsection{DÍA 3: FLUIDEZ Y PRÁCTICA}

\vspace{0.2cm}

\section{Sesión 1 (30 min): Habla sin preparación}

\vspace{0.2cm}

\textbf{Ejercicio:}

Graba 2 minutos hablando sobre "Mis vacaciones" SIN preparación.

\vspace{0.2cm}

\textbf{Luego:}

\begin{itemize}

  \item Escucha la grabación

  \item Identifica 5 áreas de mejora

  \item Repite intentando mejorar

\end{itemize}

\vspace{0.2cm}

---

\vspace{0.2cm}

\section{Sesión 2 (30 min): Presentación corta}

\vspace{0.2cm}

\textbf{Prepara y presenta:}

Tema: "Lo que hice durante las vacaciones"

\begin{itemize}

  \item Duración: 3 minutos

  \item Mínimo 5 conectores diferentes

  \item Mínimo 2 tiempos pasados diferentes

\end{itemize}

\vspace{0.2cm}

---

\vspace{0.2cm}

\section{ESTRATEGIAS DE REACTIVACIÓN}

\vspace{0.2cm}

\subsection{Si te sientes MUY oxidada/o}

\vspace{0.2cm}

1. \textbf{Revisa materiales del curso}

\begin{itemize}

  \item Lee todas las guías de sesiones 1-16

  \item Rehace ejercicios

\end{itemize}

\vspace{0.2cm}

2. \textbf{Escucha español activamente}

\begin{itemize}

  \item Podcasts, videos, música

  \item Repite en voz alta

\end{itemize}

\vspace{0.2cm}

3. \textbf{Habla contigo mismo}

\begin{itemize}

  \item Narrar tu día

  \item Describir lo que ves

  \item Practicar monólogos

\end{itemize}

\vspace{0.2cm}

\subsection{Si te sientes ALGO oxidada/o}

\vspace{0.2cm}

1. \textbf{Práctica diaria de 15 minutos}

\begin{itemize}

  \item Habla sobre tu día

  \item Lee en voz alta

\end{itemize}

\vspace{0.2cm}

2. \textbf{Revisa gramática específica}

\begin{itemize}

  \item Solo áreas problemáticas

\end{itemize}

\vspace{0.2cm}

3. \textbf{Conversación con nativo}

\begin{itemize}

  \item Intercambio de idiomas

  \item Clases particulares

\end{itemize}

\vspace{0.2cm}

\subsection{Si te sientes BIEN}

\vspace{0.2cm}

1. \textbf{Mantén práctica regular}

\begin{itemize}

  \item 30 min diarios

  \item Conversaciones

\end{itemize}

\vspace{0.2cm}

2. \textbf{Desafíate}

\begin{itemize}

  \item Temas complejos

  \item Debate

  \item Presentaciones largas

\end{itemize}

\vspace{0.2cm}

---

\vspace{0.2cm}

\section{PLAN DE VUELTA A CLASE}

\vspace{0.2cm}

\subsection{Primera semana post-vacaciones}

\vspace{0.2cm}

\textbf{Lunes:}

\begin{itemize}

  \item Repaso de conectores (30 min)

  \item Presentación: "¿Qué hice en vacaciones?" (3 min)

\end{itemize}

\vspace{0.2cm}

\textbf{Martes:}

\begin{itemize}

  \item Repaso de pretéritos (30 min)

  \item Narración: "Un momento memorable de vacaciones" (3 min)

\end{itemize}

\vspace{0.2cm}

\textbf{Miércoles:}

\begin{itemize}

  \item Repaso de subjuntivo (30 min)

  \item Opinión: "¿Debería haber vacaciones más largas?" (3 min)

\end{itemize}

\vspace{0.2cm}

\textbf{Jueves:}

\begin{itemize}

  \item Repaso de ser/estar (30 min)

  \item Descripción: "Mi lugar favorito de vacaciones" (3 min)

\end{itemize}

\vspace{0.2cm}

\textbf{Viernes:}

\begin{itemize}

  \item Repaso general (30 min)

  \item Conversación libre con compañero (5 min)

\end{itemize}

\vspace{0.2cm}

---

\vspace{0.2cm}

\section{EJERCICIOS DE REACTIVACIÓN RÁPIDA}

\vspace{0.2cm}

\subsection{Ejercicio 1: Frases rápidas}

Completa estas frases rápidamente:

\vspace{0.2cm}

1. "Si ____ (tener) dinero, ____ (viajar) por el mundo."

Respuesta: tuviera/tuviese, viajaría

\vspace{0.2cm}

2. "Ojalá ____ (haber) ____ (saber) eso antes."

Respuesta: hubiera/hubiese sabido

\vspace{0.2cm}

3. "Cuando ____ (ser) niño, ____ (ir) al parque todos los días."

Respuesta: era, iba

\vspace{0.2cm}

4. "No creo que él ____ (poder) venir."

Respuesta: pueda

\vspace{0.2cm}

5. "____ (Estar/Es) muy importante ____ (estudiar) todos los días."

Respuesta: Es, estudiar

\vspace{0.2cm}

---

\vspace{0.2cm}

\subsection{Ejercicio 2: Conectores}

Usa conectores para unir estas ideas:

\vspace{0.2cm}

\begin{itemize}

  \item Hice ejercicio

  \item Fui al supermercado

  \item Cociné

  \item Comí

\end{itemize}

\vspace{0.2cm}

\textbf{Respuesta:} "Primero hice ejercicio. Después fui al supermercado. Luego cociné. Finalmente, comí."

\vspace{0.2cm}

---

\vspace{0.2cm}

\subsection{Ejercicio 3: Mini-presentación}

Habla durante 2 minutos sobre: "Mi vuelta a la normalidad después de vacaciones"

\vspace{0.2cm}

\textbf{Incluye:}

\begin{itemize}

  \item Mínimo 5 conectores diferentes

  \item Mínimo 2 tiempos pasados

  \item Mínimo 1 expresión de opinión

\end{itemize}

\vspace{0.2cm}

---

\vspace{0.2cm}

\section{CHECKLIST DE LISTO PARA VOLVER}

\vspace{0.2cm}

\begin{itemize}

  \item [ ] He repasado ser/estar

  \item [ ] He repasado subjuntivo

  \item [ ] He repasado pretéritos

  \item [ ] He repasado condicionales

  \item [ ] Recuerdo conectores variados

  \item [ ] He practicado hablar en voz alta

  \item [ ] Me siento más confiada/o

\end{itemize}

\vspace{0.2cm}

---

\vspace{0.2cm}

\textbf{¡Bienvenido/a de vuelta! ¡Vamos a continuar tu progreso!}

\vspace{0.2cm}

---

\vspace{0.2cm}

\textbf{Sesión 17: Repaso Post-Vacaciones}

\textbf{Nivel C1 - Español Oral Avanzado}

\vspace{0.2cm}

\end{document}
